{\small
\begin{verbatim}
template <class ProblemT>
class serialLipshitzianNode : virtual public pico::branchSub
{
public:

  /// The type of the branching class
  typedef serialLipshitzian<ProblemT> branching_t;

  /// Return a pointer to the global branching object
  branching_t* global() const;

  /// Return a pointer to the base class of the global branching object
  pico::branching* bGlobal() const;

  /// Link the debugging in the subproblem to the debugging level
  /// set within the global branching object.
  REFER_DEBUG(global());

  /// Initialize a subproblem using a branching object
  void initialize(branching_t* master);

  /// Initialize a subproblem as a child of a parent subproblem.
  void initialize(serialLipshitzianNode<ProblemT>* parent,int whichChild);

  /// Initialize this subproblem to be the root of the branching tree
  void setRootComputation();

  /// Compute the lower bound on this subproblem's value.
  void boundComputation(double* controlParam);

  /// Determine how many children will be generated and how they will be
  /// generated (e.g. the branching variable).
  int splitComputation();

  /// Create a child subproblem of the current subproblem
  branchSub* makeChild(int whichChild);

  /// Returns true if this subproblem represents a feasible solution
  bool candidateSolution();

  /// Updates the incumbent solution.
  void updateIncumbent();

};
\end{verbatim}
}
